% Vietnamese translation for OI Public Library
% Source-PDF: 英文资料 ENGLISH MATERIALS/Computer Systems, A Programmer’s Perspective 3rd Edition.pdf
\documentclass[11pt]{article}
\usepackage{oipl-vn}

\title{Hệ thống máy tính: góc nhìn của lập trình viên}
\author{Randal E. Bryant, David R. O'Hallaron}
\date{Pearson, ấn bản thứ ba, 2016}

\begin{document}
\maketitle

\origtitle{Computer Systems: A Programmer's Perspective, Third Edition}
\sourcepath{English Materials / Computer Systems, A Programmer's Perspective 3rd Edition.pdf}

\renewcommand{\abstractname}{Tóm tắt}
\begin{abstract}
Đây là ghi chú dịch tiếng Việt cho sách \emph{Computer Systems: A Programmer's
Perspective}, ấn bản thứ ba. Sách giải thích hệ thống máy tính từ góc nhìn của
người lập trình: biểu diễn dữ liệu, mã máy, kiến trúc xử lý, bộ nhớ, liên kết,
ngoại lệ, bộ nhớ ảo, I/O, mạng và lập trình đồng thời. Ghi chú này chuyển ngữ
thông tin đầu sách và cấu trúc chương; các đoạn mã và bài tập lập trình không
được chuyển thành mẫu mã mới.
\end{abstract}

\section{Thông tin sách}

\begin{center}
\begin{tabular}{ll}
\textbf{Tác giả} & Randal E. Bryant; David R. O'Hallaron \\
\textbf{Cơ quan} & Carnegie Mellon University \\
\textbf{Nhà xuất bản} & Pearson \\
\textbf{Ấn bản} & Thứ ba \\
\textbf{Năm} & 2016 \\
\textbf{ISBN} & 978-0-13-409266-9
\end{tabular}
\end{center}

Sách được xây dựng từ khóa học hệ thống máy tính tại Carnegie Mellon. Điểm
nhìn trung tâm là: hiểu phần cứng, trình biên dịch, hệ điều hành và mạng sẽ
giúp lập trình viên viết chương trình nhanh hơn, đúng hơn và dễ gỡ lỗi hơn.

\section{Mục lục chương đã dịch}

\begin{description}
  \item[Chương 1. Một chuyến tham quan hệ thống máy tính]
  Giới thiệu bit, chương trình, quá trình biên dịch, bộ xử lý, bộ nhớ, cache,
  hệ điều hành, mạng, song song và các lớp trừu tượng.

  \item[Chương 2. Biểu diễn và thao tác thông tin]
  Biểu diễn số nguyên, số dấu phẩy động, phép toán bit, tràn số, mã bù hai và
  các rủi ro thường gặp khi chuyển đổi kiểu.

  \item[Chương 3. Biểu diễn chương trình ở mức máy]
  Mã hợp ngữ, thanh ghi, stack, điều khiển rẽ nhánh, vòng lặp, thủ tục, mảng,
  con trỏ, cấu trúc dữ liệu và các lỗi bộ nhớ.

  \item[Chương 4. Kiến trúc bộ xử lý]
  Mô hình tập lệnh Y86-64, thiết kế logic, thực thi tuần tự, pipeline, hazard,
  ngoại lệ và điều khiển pipeline.

  \item[Chương 5. Tối ưu hiệu năng chương trình]
  Cách trình biên dịch tối ưu, đo hiệu năng, loại bỏ tính toán thừa, unroll
  vòng lặp, tăng song song mức lệnh và hiểu giới hạn của bộ nhớ.

  \item[Chương 6. Hệ phân cấp bộ nhớ]
  Công nghệ lưu trữ, locality, cache, tổ chức cache và cách viết chương trình
  tận dụng bộ nhớ hiệu quả.

  \item[Chương 7. Liên kết]
  Liên kết tĩnh, liên kết động, ký hiệu, relocation, thư viện và cách linker
  kết hợp các module thành chương trình thực thi.

  \item[Chương 8. Luồng điều khiển ngoại lệ]
  Ngoại lệ, tiến trình, tín hiệu, chuyển ngữ cảnh, shell, fork, exec và wait.

  \item[Chương 9. Bộ nhớ ảo]
  Không gian địa chỉ, phân trang, dịch địa chỉ, cấp phát động, garbage
  collection và các lỗi bộ nhớ phổ biến.

  \item[Chương 10. I/O mức hệ thống]
  File, descriptor, I/O Unix, metadata, redirection và các hàm đọc ghi tin cậy.

  \item[Chương 11. Lập trình mạng]
  Socket, mô hình client-server, địa chỉ Internet, DNS, HTTP và xây dựng dịch
  vụ mạng đơn giản.

  \item[Chương 12. Lập trình đồng thời]
  Thread, semaphore, mutex, condition variable, race, deadlock và các mô hình
  server đồng thời.
\end{description}

\section{Ghi chú nội dung}

Trong bối cảnh OI, nhiều phần của sách nằm ngoài phạm vi thi thuật toán trực
tiếp, nhưng vẫn hữu ích cho lập trình hiệu năng cao và hiểu lỗi hệ thống:
\begin{itemize}
  \item Chương 2--3 giúp hiểu biểu diễn số, overflow, bit trick và layout bộ nhớ.
  \item Chương 5--6 liên quan trực tiếp tới tối ưu hằng số, cache và locality.
  \item Chương 7--10 hỗ trợ gỡ lỗi chương trình lớn, lỗi liên kết, lỗi bộ nhớ
        và thao tác file.
  \item Chương 11--12 hữu ích cho hệ thống, server và lập trình đồng thời, nhưng
        không phải trọng tâm OI cổ điển.
\end{itemize}

\section{Ghi chú về trích xuất}

Tệp PDF có 1120 trang và được xử lý qua OCR của Acrobat. Văn bản tiếng Anh
trích xuất được nhưng có lỗi nhận dạng ở tiêu đề, dấu câu và một số ký hiệu.
Vì sách chứa nhiều đoạn mã, hình và bảng, bản dịch đầy đủ cần đối chiếu trang
gốc khi chuyển LaTeX; các đoạn mã nên được giữ như mã tham khảo thay vì dịch
thành mẫu code mới.

\end{document}
